\begin{statementbox}
	\textbf{\textcolor{CStatement}{条件}}
	\begin{description}[style=nextline, leftmargin=2.8em, labelsep=0.5em, itemsep=0.35em]
		\item[\textbf{\textcolor{CStatement}{条件1（两圆交于两点）：}}] 两个不同圆 $C_1,C_2$ 相交于两个不同的点 $A,B$，圆心分别为 $O_1(a_1,b_1)$、$O_2(a_2,b_2)$，半径分别为 $r_1,r_2>0$，圆心距 $d=|O_1O_2|>0$。等价地，$|r_1-r_2|<d<r_1+r_2$。
		\item[\textbf{\textcolor{CStatement}{条件2（方程表示）：}}] 两圆方程分别为
		\[
			C_i:\quad (x-a_i)^2+(y-b_i)^2=r_i^2\quad(i=1,2).
		\]
	\end{description}

	\textbf{\textcolor{CStatement}{结论}}
	\begin{description}[style=nextline, leftmargin=2.8em, labelsep=0.5em, itemsep=0.45em]
		\item[\textbf{\textcolor{CStatement}{结论一（几何直觉：方向加中点）：}}]
		\textbf{\textcolor{CStatement}{直观描述：}} 公共弦 $AB$ 垂直于连心线 $O_1O_2$，并在点 $M$ 处被连心线平分。以 $O_1\to O_2$ 为正方向，设 $t=\overline{O_1M}$ 为有向距离，则
		\[
			\boxed{t=\frac{d^2+r_1^2-r_2^2}{2d}},\qquad
			\boxed{M=O_1+\frac{t}{d}(O_2-O_1)}.
		\]
		于是公共弦所在直线可写成
		\[
			\boxed{(a_2-a_1)(x-x_M)+(b_2-b_1)(y-y_M)=0}.
		\]

		\item[\textbf{\textcolor{CStatement}{结论二（代数捷径：两圆方程相减）：}}]
		\textbf{\textcolor{CStatement}{直观描述：}} 共同交点同时满足两个圆方程；将方程左端相减，$x^2+y^2$ 消去，所得一次方程正是公共弦所在直线：
		\[
			\boxed{2(a_2-a_1)(x-a_1)+2(b_2-b_1)(y-b_1)
			=d^2+r_1^2-r_2^2}.
		\]

		\item[\textbf{\textcolor{CStatement}{推广写法（一般式）：}}]
		若两圆为 $x^2+y^2+D_ix+E_iy+F_i=0$（$i=1,2$），则
		\[
			\boxed{(D_1-D_2)x+(E_1-E_2)y+F_1-F_2=0}.
		\]
	\end{description}
\end{statementbox}
